\chapter{Nonlinearity: the Bratu problem and Newton's method}
\label{ch:B1_bratu_newton}
\noindent\texttt{tutorials/B\_nonlinear/B1\_bratu\_newton.py}\\[0.75em]
\begin{outcome}
You can take a solver you trust for LINEAR problems and
extend it to a nonlinear PDE with a hand-written Newton loop: residual,
Jacobian, solve, update — and you can PROVE it works two independent ways:
quadratic convergence of the Newton iteration, and unchanged MMS spatial
orders (2 for p1, 3 for p2). Nonlinearity changes the iteration, not the
discretization accuracy.
\end{outcome}

The Bratu (solid-fuel ignition) problem:
    -lap(u) = lam * exp(u) + f,     u = g on the boundary.
Weak residual and Jacobian at the current iterate u\_h:
    R\_a(u) = (grad N\_a, grad u\_h) - (N\_a, lam e\textasciicircum{}\{u\_h\} + f)
    J\_ab   = (grad N\_a, grad N\_b) - (N\_a, lam e\textasciicircum{}\{u\_h\} N\_b)
J is the stiffness matrix MINUS an exp-weighted mass matrix — we assemble
that weighted mass right here in \textasciitilde{}10 lines of numpy einsum over element
batches (a pattern you will reuse constantly). MMS: pick u*, set
f = -lap(u*) - lam e\textasciicircum{}\{u*\}.

\begin{expected}
\begin{verbatim}
Newton residuals (level 5, p1, lam=3): each step roughly SQUARES the
    previous residual norm — e.g. 1e-1 -> 1e-3 -> 1e-7 -> 1e-14 (4 steps).
    Spatial orders: p1 -> ~2.0, p2 -> ~3.0, exactly as in A1.
\end{verbatim}
\end{expected}

\section*{The script}
\begin{lstlisting}
import numpy as np
import scipy.sparse as sp
from scipy.sparse.linalg import splu

from diffsim.octree.build import build_uniform
from diffsim.mesh.nodes import build_mesh
from diffsim.mesh.constraints import build_constraints
from diffsim.mesh.basis import basis_tables
from diffsim.assembly.operators import DeviceMesh, assemble_csr
from diffsim.physics.poisson import gauss_points, l2_error_masked

DEVICE = "cuda:0"
LAM = 3.0
u_star = lambda x: np.sin(np.pi * x[:, 0]) * np.sin(np.pi * x[:, 1])
f_star = lambda x: 2 * np.pi ** 2 * u_star(x) - LAM * np.exp(u_star(x))


class BratuWorkspace:
    """Everything reusable across Newton steps for one mesh."""

    def __init__(self, level, p):
        tree = build_uniform(level, dim=2)
        self.mesh = build_mesh(tree, p=p)
        self.cons = build_constraints(self.mesh)
        tables = basis_tables(p, dim=2)
        self.dm = DeviceMesh.from_mesh(self.mesh, self.cons, tables, DEVICE)
        self.K = assemble_csr(self.dm)            # linear part, assembled once
        self.T = self.cons.T.tocsr()
        self.xq = gauss_points(self.mesh, self.dm.tables_by_p)
        coords = self.mesh.node_coords[self.cons.free_nodes]
        self.bdry = np.where(self.mesh.boundary_nodes[self.cons.free_nodes])[0]
        self.gvals = u_star(coords[self.bdry])
        self.n_free = self.T.shape[1]

    def gp_values(self, u_free):
        """u_h at every Gauss point, per p-bin (the einsum you will reuse)."""
        full = np.asarray(self.T @ u_free)
        out = {}
        for pv, b in self.dm.bins.items():
            tb = self.dm.tables_by_p[pv]
            out[pv] = np.einsum("qa,ea->eq", tb.N,
                                full[self.mesh.conn_of[pv]])
        return out

    def weighted_integrals(self, w_gp):
        """(vector (N_a, w),  matrix (N_a, w N_b)) for a GP field w."""
        dm, mesh = self.dm, self.mesh
        Fv = np.zeros(dm.n_nodes)
        rows, cols, vals = [], [], []
        for pv, b in dm.bins.items():
            tb = dm.tables_by_p[pv]
            h = mesh.tree.h()[mesh.bins[pv]]
            jac = (h / 2.0) ** dm.dim
            conn = mesh.conn_of[pv].astype(np.int64)
            wq = w_gp[pv]
            be = np.einsum("qa,eq,q,e->ea", tb.N, wq, tb.w, jac)
            np.add.at(Fv, conn.ravel(), be.ravel())
            Me = np.einsum("qa,qb,eq,q,e->eab", tb.N, tb.N, wq, tb.w, jac)
            nbf = conn.shape[1]
            rows.append(np.repeat(conn, nbf, axis=1).ravel())
            cols.append(np.tile(conn, (1, nbf)).ravel())
            vals.append(Me.ravel())
        M = sp.coo_matrix((np.concatenate(vals),
                           (np.concatenate(rows), np.concatenate(cols))),
                          shape=(dm.n_nodes,) * 2).tocsr()
        return (np.asarray(self.T.T @ Fv),
                (self.T.T @ M @ self.T).tocsr())


def newton_solve(ws, tol=1e-12, verbose=False):
    u = np.zeros(ws.n_free)
    u[ws.bdry] = ws.gvals                        # start feasible on boundary
    hist = []
    for it in range(20):
        uq = ws.gp_values(u)
        w_exp = {pv: LAM * np.exp(uq[pv]) for pv in uq}          # lam e^u
        w_rhs = {pv: w_exp[pv] + f_star(ws.xq[pv]).reshape(
            uq[pv].shape) for pv in uq}                          # + f
        b_nl, M_exp = ws.weighted_integrals(w_rhs)
        _, M_jac = ws.weighted_integrals(w_exp)
        R = ws.K @ u - b_nl                      # residual
        J = (ws.K - M_jac).tolil()               # Jacobian
        R[ws.bdry] = u[ws.bdry] - ws.gvals       # Dirichlet rows
        for i in ws.bdry:
            J.rows[i] = [int(i)]
            J.data[i] = [1.0]
        rn = np.linalg.norm(R, np.inf)
        hist.append(rn)
        if rn < tol:
            break
        u = u - splu(J.tocsr().tocsc()).solve(R)
    if verbose:
        print("   Newton |R|_inf per step: "
              + "  ".join(f"{r:.2e}" for r in hist))
    return u


if __name__ == "__main__":
    ws = BratuWorkspace(5, 1)
    newton_solve(ws, verbose=True)               # watch it square
    for p, levels in ((1, (4, 5, 6)), (2, (3, 4, 5))):
        errs = []
        for lv in levels:
            ws = BratuWorkspace(lv, p)
            u = newton_solve(ws)
            u_all = np.asarray(ws.T @ u)
            errs.append(l2_error_masked(ws.dm, u_all, u_star,
                                        lambda x: np.ones(len(x), bool)))
        orders = [np.log2(errs[i] / errs[i + 1]) for i in range(2)]
        print(f"p={p}: errors " + "  ".join(f"{e:.3e}" for e in errs)
              + "   orders " + "  ".join(f"{o:.2f}" for o in orders))
\end{lstlisting}

\begin{explore}
\begin{verbatim}
(a) Replace the Jacobian by the FROZEN one from the first iterate
      (modified Newton). The iteration still converges — at what rate?
      Count iterations to 1e-12 for both variants.
  (b) Drop the -M_jac term entirely (Picard/fixed-point). Rate now?
  (c) The physical Bratu problem (f = 0, g = 0) has TWO solutions for
      lam below a critical lam* ~ 6.81 (2-D square) and none above. Track
      the lower branch: solve at lam = 1, 2, ..., using each solution as
      the next initial guess. What happens to the Newton step count as
      lam -> lam*? (B2 turns this into proper continuation.)
  (d) PERFORMANCE CORNER: this Newton loop re-assembles the weighted mass
      on the HOST every step. Time host-assembly vs the splu solve per
      Newton step at level 7. Which dominates? Sketch (do not implement)
      what a warp kernel for the exp-weighted mass would look like — the
      brick API of the README is exactly that hole.
\end{verbatim}
\end{explore}
